Neural interfaces in biohybrid systems use stimulation and recording to interact with living networks. Under a limited recording budget, choosing where to record becomes part of predicting a specified stimulus response. Goal-oriented design ranks recording sets by their estimated prediction error. We test whether averaging local risk over a parameter prior improves the resulting response prediction or the accuracy of the risk estimate itself.

We compared single-response and shared three-response objectives, each using a nominal-parameter or prior-averaged risk criterion, in simulated 12-unit networks with known connectivity and uncertain neuronal gains and biases. The target was the time-integrated difference between stimulated and control activity. All methods used the same feasible recording sites, measurement budget and maximum a posteriori estimator. The primary evaluation used 40 independent networks under broad, correctly specified parameter uncertainty, with three recorded neurons and stimulation and readout sites excluded. Each network was simulated with several parameter draws, noise draws and readouts; networks were the unit of analysis.

Compared with the nominal criterion, prior averaging changed aggregate mean squared response error by -3.74\% for the single-response objective (97.5\% paired network-cluster confidence interval, -18.48\% to 12.12\%) and +2.69\% for the three-response objective (-2.95\% to 11.51\%). Both confidence intervals included zero. The two criteria selected identical recording sets for 82 of 120 target-specific designs (three per network) and 26 of 40 shared designs. To evaluate risk estimates independently of recording selection, a separately specified study used 40 new networks with recordings, data and point predictions held fixed. The primary ensemble-minus-nominal score difference was -0.45565 (95\% confidence interval, -1.15857 to 0.02367), using log(v) + e\textsuperscript{2}/v, where v is the estimated mean squared error and e is the response prediction error. Nineteen of 40 networks favoured the ensemble forecast, one network supplied about 64\% of the summed difference, and the nominal forecast risk was closer to the actual error in aggregate.

In an exploratory extension to eight independent 100-unit rank-two networks, four of eight settings had 95\% intervals that excluded zero, all favouring prior averaging; none excluded zero in the other direction, and one setting reversed in sign. Under the shared fitting procedure, prior averaging can affect response prediction only by changing which neurons are recorded.
